\NSPage{055}%
\begin{EESegment}{EE-TGT-P055-001}
For a tuple \NSi{NS-F-P055-001}{V} of physical fields, write
\NSFormula{NS-F-P055-002}{display}{5.23}
\begin{equation}
 |V|_m=\max_{|\alpha|+b\leq m}|\partial_x^\alpha\partial_t^bV|.
 \tag{5.23}\label{eq:5.23}
\end{equation}
At the point in question, this is the largest absolute value among all Cartesian space--time derivatives of the tuple
through order \NSi{NS-F-P055-003}{m}. Write
\NSi{NS-F-P055-004}{u_n=u_{r,n}e_r+u_{\theta,n}e_\theta+u_{z,n}e_z} for the physical velocity coefficient in
\eqref{eq:5.1}, and write \NSi{NS-F-P055-005}{\mathcal T_0} for the leading stress profile from \eqref{eq:4.11}.
\end{EESegment}

\begin{EESegment}{EE-TGT-P055-002}
Before putting any cutoff in \NSi{NS-F-P055-006}{q}, define the finite tuple
\NSFormula{NS-F-P055-007}{display}{none}
\[
 \mathcal U^{[N]}=
 (u_{\mathrm{slow}}^{[N]},p_{\mathrm{slow}}^{[N]},T_{\mathrm{slow}}^{[N]})
 =\sum_{n=0}^{N}(u_n,p_n,q^{-A-1/2+\lambda_n}\mathcal T_n).
\]
The profiles in this sum already contain their radial cutoffs and moment corrections. For a physical stress pair
\NSi{NS-F-P055-008}{T=(T_\theta,T_z)}, define
\NSFormula{NS-F-P055-009}{display}{5.24}
\begin{equation}
 \mathcal F_{\mathrm{slow}}(u,p,T)=\mathscr R(u,p)
 +(\partial_r+2/r)T_\theta e_\theta+(\partial_r+1/r)T_ze_z.
 \tag{5.24}\label{eq:5.24}
\end{equation}
The last two terms are the tangential divergence of the stress that is being kept for the later wave construction.
\end{EESegment}

\begin{proposition}\label{prop:5.3}
\begin{EESegment}{EE-TGT-P055-003}
Fix the leading profiles from Theorem~\ref{thm:4.6}. There is a sequence of positive-order profiles, with each order
built first by Lemma~\ref{lem:5.1} and then by Lemma~\ref{lem:5.2} after all lower orders have been finalized. Every
finite velocity sum \NSi{NS-F-P055-010}{u_{\mathrm{slow}}^{[N]}} is divergence-free. Take any fixed compact profile
range \NSi{NS-F-P055-011}{0\leq X\leq X_{\max}}, \NSi{NS-F-P055-012}{-1\leq\eta\leq1}, with
\NSi{NS-F-P055-013}{0<X_{\max}<\infty}. For every \NSi{NS-F-P055-014}{N,m\geq0} and
\NSi{NS-F-P055-015}{0<q\leq1}, the augmented momentum residual obeys
\NSFormula{NS-F-P055-016}{display}{5.25}
\begin{equation}
 |\mathcal F_{\mathrm{slow}}(\mathcal U^{[N]})|_m
 \leq C_{N,m}q^{2h(N+1)-K_m},
 \qquad K_m\ \text{independent of }N.
 \tag{5.25}\label{eq:5.25}
\end{equation}
The constants can depend on the fixed profile range. The constant \NSi{NS-F-P055-017}{C_{N,m}} can depend on
\NSi{NS-F-P055-018}{N}, while the loss \NSi{NS-F-P055-019}{K_m} in the exponent does not.
\end{EESegment}
\end{proposition}

\begin{proof}
\begin{EESegment}{EE-TGT-P055-004}
\textbf{Step 1. Build the sequence and cancel every coefficient that has been kept.}
At order \NSi{NS-F-P055-020}{n}, apply Lemma~\ref{lem:5.1} with the finalized orders
\NSi{NS-F-P055-021}{0,\ldots,n-1}. Next apply Lemma~\ref{lem:5.2}. This extends and corrects the current profiles,
rebuilds \NSi{NS-F-P055-022}{V_n,\Pi_n} from Equations~\eqref{eq:5.2} and~\eqref{eq:5.5}, and defines
\NSi{NS-F-P055-023}{\mathcal T_n} by \eqref{eq:5.9}. These operations keep the equations unchanged on the inner
interval that was retained, and they keep the regularity needed to form \NSi{NS-F-P055-024}{\Omega_n/X} at the next
order. Both lemmas use the same fixed radial intervals, so this process continues for every
\NSi{NS-F-P055-025}{n}.
\end{EESegment}

\begin{EESegment}{EE-TGT-P055-005}
Equation~\eqref{eq:5.2} gives exact incompressibility at each order, and hence for every finite sum. The pressure equation
holds everywhere. Differentiating the stress primitives gives
\NSFormula{NS-F-P055-026}{display}{none}
\[
 (\partial_R+2/R)\mathcal T_{n,\theta}=-\mathfrak r_{\theta,n},
 \qquad
 (\partial_R+1/R)\mathcal T_{n,z}=-\mathfrak r_{z,n}.
\]
Thus the tangential residual at each retained order is exactly the negative of the stress divergence. Proposition~\ref{prop:4.2}
gives the same leading identities at order zero. Once the sum has been cut off at \NSi{NS-F-P055-027}{N}, every product
or shifted viscosity term that has not been canceled has an exponent increment of at least
\NSi{NS-F-P055-028}{2h(N+1)}. For example, when \NSi{NS-F-P055-029}{N=1}, both the product of two order-one velocity
coefficients and axial viscosity applied to order one first appear at order two.
\end{EESegment}

\begin{EESegment}{EE-TGT-P055-006}
\textbf{Step 2. Estimate the physical derivatives of the finite residual.}
The coordinate formulas show exactly how physical differentiation changes the powers of \NSi{NS-F-P055-030}{q}.
Suppose that \NSi{NS-F-P055-031}{g(x_\perp/\sqrt q,\eta)} is smooth in its Cartesian arguments on a compact profile
range. If there are \NSi{NS-F-P055-032}{a=|\beta|} transverse derivatives, \NSi{NS-F-P055-033}{k} axial derivatives,
and \NSi{NS-F-P055-034}{l} time derivatives, then
\NSFormula{NS-F-P055-035}{display}{5.26}
\begin{equation}
 \left|\partial_{x_\perp}^{\beta}\partial_z^k\partial_t^l
 \bigl[q^b g(x_\perp/\sqrt q,\eta)\bigr]\right|
 \leq C_m(1+|b|)^m\|g\|_{C^m}q^{b-a/2-Dk-l},
 \qquad m=a+k+l.
 \tag{5.26}\label{eq:5.26}
\end{equation}
\end{EESegment}

\NSPage{056}%
\begin{EESegment}{EE-TGT-P056-001}
This follows from the chain rule and the identities
\NSi{NS-F-P056-001}{\partial_t=q^{-1}L^{-1}(-q\partial_q+X\partial_X+D\eta\partial_\eta)} and
\NSi{NS-F-P056-002}{\partial_z=q^{-D}L^{-1}(2\eta q\partial_q-2\eta X\partial_X+d\partial_\eta)}. Each transverse
derivative of these smooth Cartesian profiles contributes a factor \NSi{NS-F-P056-003}{q^{-1/2}}.
\end{EESegment}

\begin{EESegment}{EE-TGT-P056-002}
The powers of \NSi{NS-F-P056-004}{n} produced by differentiating \NSi{NS-F-P056-005}{q^{\lambda_n}} are
absorbed into the coefficient constants. The reduction in the exponent depends only on how many physical derivatives of
each kind are taken. On the common stress support, \NSi{NS-F-P056-006}{X\geq X_a>0}, so
\NSi{NS-F-P056-007}{r^{-1}=q^{-1/2}(2X)^{-1/2}} also has derivative losses that do not depend on
\NSi{NS-F-P056-008}{N}. Near the axis, use the smooth Cartesian representative of the velocity instead. A fixed
truncated residual contains only finitely many terms. Applying \eqref{eq:5.26} to those terms proves \eqref{eq:5.25},
with \NSi{NS-F-P056-009}{K_m} independent of \NSi{NS-F-P056-010}{N}.
\end{EESegment}
\end{proof}

\begin{EESegment}{EE-TGT-P056-003}
\subsection{Adding the coefficients while keeping zero divergence}\label{sec:5.4}
The finite tuples \NSi{NS-F-P056-011}{\mathcal U^{[N]}} have residuals with better and better decay in
\eqref{eq:5.25}, but the constants \NSi{NS-F-P056-012}{C_{n,m}} in the coefficient bounds \eqref{eq:5.17} can grow too
quickly for the entire formal series \eqref{eq:5.1} to converge. To obtain smooth fields, multiply the
\NSi{NS-F-P056-013}{n}\,th positive-order term by \NSi{NS-F-P056-014}{\chi(c_nq)}. Here
\NSi{NS-F-P056-015}{\chi=1} near zero and \NSi{NS-F-P056-016}{\chi=0} once its argument reaches one, and the constants are chosen with
\NSi{NS-F-P056-017}{c_n\to\infty}, as in \eqref{eq:5.34}. Each later term is then zero outside a successively smaller neighborhood
of \NSi{NS-F-P056-018}{q=0}. On any compact set where \NSi{NS-F-P056-019}{q} stays a positive distance away from zero,
only finitely many terms remain. The remaining point is to show that the finite residual bound in
Proposition~\ref{prop:5.3} survives these cutoffs. Lemma~\ref{lem:5.4} proves this in \eqref{eq:5.36} for the augmented
residual \eqref{eq:5.24}.
\end{EESegment}

\begin{EESegment}{EE-TGT-P056-004}
To keep the velocity incompressible, apply the cutoffs to vector potentials before differentiating them. An axisymmetric
Stokes streamfunction \NSi{NS-F-P056-020}{S} gives such a potential. If \NSi{NS-F-P056-021}{S/r^2} is smooth as a
function of \NSi{NS-F-P056-022}{(r^2,z,t)}, then
\NSFormula{NS-F-P056-023}{display}{none}
\[
 \mathcal A=\frac Sr e_\theta=\frac S{r^2}(-x_2,x_1,0),
 \qquad
 \operatorname{curl}\mathcal A=-\frac{\partial_zS}{r}e_r+\frac{\partial_rS}{r}e_z.
\]
The smoothness of \NSi{NS-F-P056-024}{S/r^2} in \NSi{NS-F-P056-025}{(r^2,z,t)} makes this potential smooth in
Cartesian coordinates.
\end{EESegment}

\begin{EESegment}{EE-TGT-P056-005}
For the positive-order profiles already constructed, the physical streamfunctions and vector potentials are
\NSFormula{NS-F-P056-026}{display}{5.27}
\begin{equation}
 F_n=X\mathcal A_X(U_n),
 \qquad \mathcal S_n=q^{1-A+\lambda_n}F_n,
 \qquad \mathcal A_n=(\mathcal S_n/r)e_\theta,
 \qquad
 u_{z,n}=\partial_s\mathcal S_n,
 \qquad ru_{r,n}=-\partial_z\mathcal S_n.
 \tag{5.27}\label{eq:5.27}
\end{equation}
Here \NSi{NS-F-P056-027}{s=r^2/2}, as in \eqref{eq:4.1}. In fact,
\NSi{NS-F-P056-028}{\operatorname{curl}\mathcal A_n=u_{r,n}e_r+u_{z,n}e_z}, so the curl gives exactly the radial and
axial velocity components. In Cartesian coordinates, \NSi{NS-F-P056-029}{r^2=2qX} gives
\NSFormula{NS-F-P056-030}{display}{none}
\[
 \mathcal A_n=\frac{\mathcal S_n}{r^2}(-x_2,x_1,0)
 =\frac12q^{-A+\lambda_n}\frac{F_n}{X}(-x_2,x_1,0).
\]
Since \NSi{NS-F-P056-031}{F_n/X} is smooth at \NSi{NS-F-P056-032}{X=0}, this formula extends smoothly through the
axis whenever \NSi{NS-F-P056-033}{q>0}. Lemma~\ref{lem:5.2} gives common supports and bounds for every fixed derivative
of each coefficient. No bound uniform in \NSi{NS-F-P056-034}{n} is needed.
\end{EESegment}

\begin{EESegment}{EE-TGT-P056-006}
Cutting off \NSi{NS-F-P056-035}{\mathcal A_n} and then taking its curl includes the extra velocity term that comes from
differentiating the cutoff, so the result remains divergence-free. The summation argument will be stated for tuples of
physical fields, which also allows it to be used later for the wave and mean correction sequence. In the present case,
the index \NSi{NS-F-P056-036}{j=n}, the decay order is \NSi{NS-F-P056-037}{g_n=2nh}, and the entries are
\NSFormula{NS-F-P056-038}{display}{none}
\[
 A_n=\mathcal A_n,
 \qquad B_n=u_{\theta,n}e_\theta,
 \qquad p_n=q^{-2A+\lambda_n}\Pi_n,
 \qquad T_n=q^{-A-1/2+\lambda_n}\mathcal T_n.
\]
The differential polynomial is \NSi{NS-F-P056-039}{\mathcal F_{\mathrm{slow}}} from \eqref{eq:5.24}. Inside the lemma,
\NSi{NS-F-P056-040}{U_j} names a tuple of physical fields, not the axial profile
\NSi{NS-F-P056-041}{U_n(X,\eta)}. Use the derivative notation in \eqref{eq:5.23}, and set
\NSi{NS-F-P056-042}{\Lambda_{\log}(q)=1+|\log q|}. Every estimate in the lemma is uniform on the domain stated there.
\end{EESegment}

\NSPage{057}%
\begin{EESegment}{EE-TGT-P057-001}
The general argument needs one common domain, decay orders that rise to infinity, and derivative losses that do not depend
on the summation index. Let \NSi{NS-F-P057-001}{\Omega} be a domain on which
\NSi{NS-F-P057-002}{0<q<q_0\leq1}. Suppose \NSi{NS-F-P057-003}{q=q(z,t)} is smooth, stays a positive distance away
from zero on every compact subset of \NSi{NS-F-P057-004}{\Omega}, and satisfies
\NSFormula{NS-F-P057-005}{display}{5.28}
\begin{equation}
 |\partial_z^a\partial_t^bq|\leq C_{a,b}q^{1-aD-b},
 \qquad a,b\geq0,
 \qquad 0<D<1.
 \tag{5.28}\label{eq:5.28}
\end{equation}
\end{EESegment}

\begin{EESegment}{EE-TGT-P057-002}
Let \NSi{NS-F-P057-006}{U_0=(u_0,p_0,T_0)} be smooth, with
\NSi{NS-F-P057-007}{\operatorname{div}u_0=0}. Suppose that for every \NSi{NS-F-P057-008}{m}, there are finite
constants \NSi{NS-F-P057-009}{C_m,K_m,P_m} such that
\NSi{NS-F-P057-010}{|U_0|_m\leq C_mq^{-K_m}\Lambda_{\log}(q)^{P_m}}. The stress entries
\NSi{NS-F-P057-011}{T_0} and \NSi{NS-F-P057-012}{T_j} below can be left out. For
\NSi{NS-F-P057-013}{j\geq1}, suppose the smooth physical tuples
\NSFormula{NS-F-P057-014}{display}{5.29}
\begin{equation}
 Z_j=(A_j,B_j,p_j,T_j),
 \qquad B_j=b_j(r,z,t)e_\theta,
 \qquad U_j=(\operatorname{curl}A_j+B_j,p_j,T_j)
 \tag{5.29}\label{eq:5.29}
\end{equation}
are all defined on this same domain. Their Cartesian representatives must be smooth at the axis, and extending them by
zero across any side boundary of their spatial support must also be smooth.
\end{EESegment}

\begin{EESegment}{EE-TGT-P057-003}
Suppose that for every \NSi{NS-F-P057-015}{j,m}, there are finite constants
\NSi{NS-F-P057-016}{C_{j,m},P_{j,m}} such that
\NSFormula{NS-F-P057-017}{display}{5.30}
\begin{equation}
 |Z_j|_m\leq C_{j,m}q^{g_j-\ell_m}\Lambda_{\log}(q)^{P_{j,m}},
 \qquad 0<g_1\leq g_2\leq\cdots\longrightarrow\infty,
 \tag{5.30}\label{eq:5.30}
\end{equation}
where the losses \NSi{NS-F-P057-018}{\ell_m} are nonnegative, do not decrease as the derivative order rises, and do not depend on
\NSi{NS-F-P057-019}{j}.
\end{EESegment}

\begin{EESegment}{EE-TGT-P057-004}
For each finite \NSi{NS-F-P057-020}{J}, set
\NSi{NS-F-P057-021}{U^{[J]}=U_0+\sum_{j=1}^{J}U_j}. Write the scalar components of a tuple as
\NSi{NS-F-P057-022}{U=(U^1,\ldots,U^{N_U})}. A differential polynomial of fixed finite order and degree is a finite
sum of products of field components and their derivatives. Component by component, it has the form
\NSFormula{NS-F-P057-023}{display}{5.31}
\begin{equation}
 (\mathcal F(U))^a=c_{a0}(x,t)
 +\sum_{\nu=1}^{M_a}c_{a\nu}(x,t)
 \prod_{r=1}^{d_{a\nu}}\partial_{x,t}^{\beta_{a\nu r}}U^{k_{a\nu r}}.
 \tag{5.31}\label{eq:5.31}
\end{equation}
The output index ranges over the fixed finite set \NSi{NS-F-P057-024}{1\leq a\leq N_F}, and every
\NSi{NS-F-P057-025}{M_a} is finite. Fixed integers \NSi{NS-F-P057-026}{s\geq0,d\geq1} bound the derivative order and
the degree, with \NSi{NS-F-P057-027}{1\leq d_{a\nu}\leq d,\ |\beta_{a\nu r}|\leq s} and
\NSi{NS-F-P057-028}{1\leq k_{a\nu r}\leq N_U}. For
\NSi{NS-F-P057-029}{\beta=(\alpha,b)\in\mathbb N_0^4}, the notation means
\NSi{NS-F-P057-030}{\partial_{x,t}^{\beta}=\partial_x^\alpha\partial_t^b}. The coefficient functions, the indices,
and the bounds \NSi{NS-F-P057-031}{s,d} are all fixed independently of the summation and truncation indices.
\end{EESegment}

\begin{EESegment}{EE-TGT-P057-005}
For every nonconstant product in \eqref{eq:5.31}, choose one region
\NSi{NS-F-P057-032}{\Omega_{a\nu}} that contains the supports of all its field factors, both in the base and in every
increment, before and after the cutoffs are applied. For the term that does not depend on
\NSi{NS-F-P057-033}{U}, set \NSi{NS-F-P057-034}{\Omega_{a0}=\Omega}. Every coefficient
\NSi{NS-F-P057-035}{c_{a\nu}}, with \NSi{NS-F-P057-036}{0\leq\nu\leq M_a}, must satisfy the following bound for
every Cartesian space--time multi-index \NSi{NS-F-P057-037}{I\in\mathbb N_0^4}.
\NSFormula{NS-F-P057-038}{display}{5.32}
\begin{equation}
 |\partial_{x,t}^{I}c_{a\nu}(x,t)|
 \leq C_Iq^{-K_I}\Lambda_{\log}(q)^{P_I}
 \qquad ((x,t)\in\Omega_{a\nu}).
 \tag{5.32}\label{eq:5.32}
\end{equation}
The regions do not depend on the truncation. Since there are only finitely many coefficients, the finite constants
\NSi{NS-F-P057-039}{C_I,K_I,P_I} can be chosen, for each fixed \NSi{NS-F-P057-040}{I}, to work for the whole list.
They do not depend on the truncation index. The bounds therefore hold throughout the comparison between the partial sums
and the cutoff sums. The terms \NSi{NS-F-P057-041}{c_{a0}}, which do not depend on
\NSi{NS-F-P057-042}{U}, cancel in that comparison.
\end{EESegment}

\begin{EESegment}{EE-TGT-P057-006}
The next assumption concerns the full \NSi{NS-F-P057-043}{\mathcal F}, including the terms that do not depend on the
fields. Suppose every finite partial sum satisfies
\NSFormula{NS-F-P057-044}{display}{5.33}
\begin{equation}
 |\mathcal F(U^{[J]})|_m
 \leq C_{J,m}q^{\rho_J-K_m^{\mathcal F}}\Lambda_{\log}(q)^{P_{J,m}}+E_{J,m},
 \qquad \rho_J\longrightarrow\infty,
 \tag{5.33}\label{eq:5.33}
\end{equation}
where \NSi{NS-F-P057-045}{K_m^{\mathcal F}} does not depend on \NSi{NS-F-P057-046}{J}, and, for each fixed
\NSi{NS-F-P057-047}{J,m,N},
\NSi{NS-F-P057-048}{0\leq E_{J,m}\leq C_{J,m,N}q^N}. At that fixed stage, all of these bounds, including the flat
remainder, are uniform on the entire domain.
\end{EESegment}

\begin{lemma}\label{lem:5.4}
\begin{EESegment}{EE-TGT-P057-007}
There are numbers \NSi{NS-F-P057-049}{a_{j+1}\geq2a_j}, with
\NSi{NS-F-P057-050}{a_1^{-1}<q_0}, and one fixed smooth cutoff \NSi{NS-F-P057-051}{\chi} that equals one on
\NSi{NS-F-P057-052}{[0,1/2]} and zero on \NSi{NS-F-P057-053}{[1,\infty)}, such that
\NSFormula{NS-F-P057-054}{display}{5.34}
\begin{equation}
 U=U_0+\sum_{j\geq1}
 \bigl(\operatorname{curl}(\chi(a_jq)A_j)+\chi(a_jq)B_j,
       \chi(a_jq)p_j,\chi(a_jq)T_j\bigr)
 \tag{5.34}\label{eq:5.34}
\end{equation}
\end{EESegment}

\NSPage{058}%
\begin{EESegment}{EE-TGT-P058-001}
is locally finite and smooth, and its velocity component \NSi{NS-F-P058-001}{u} satisfies
\NSi{NS-F-P058-002}{\operatorname{div}u=0}. Every added term extends smoothly by zero across the outer boundary
\NSi{NS-F-P058-003}{q=q_0}. For every \NSi{NS-F-P058-004}{m} and every
\NSi{NS-F-P058-005}{J\geq\max(1,m)},
\NSFormula{NS-F-P058-006}{display}{5.35}
\begin{equation}
 |U-U^{[J]}|_m\leq2^{-J}q^{g_{J+1}/2-\ell_m'}
 \qquad\left(0<q<\frac1{2a_J}\right),
 \tag{5.35}\label{eq:5.35}
\end{equation}
where \NSi{NS-F-P058-007}{\ell_m'} does not depend on \NSi{NS-F-P058-008}{J}.

Under the residual assumption \eqref{eq:5.33}, every Cartesian space--time derivative of
\NSi{NS-F-P058-009}{\mathcal F(U)} vanishes to infinite order as \NSi{NS-F-P058-010}{q\downarrow0}. In other words,
this says
\NSFormula{NS-F-P058-011}{display}{5.36}
\begin{equation}
 \forall m,N\qquad\exists\delta,C>0\qquad
 |\mathcal F(U)|_m\leq Cq^N
 \qquad(0<q<\delta).
 \tag{5.36}\label{eq:5.36}
\end{equation}
\end{EESegment}
\end{lemma}

\begin{proof}
\begin{EESegment}{EE-TGT-P058-002}
First bound the derivatives of every cutoff. Then choose the cutoffs so that the tail remains summable after any fixed
number of derivatives, and compare the residual with one finite partial sum before the cutoffs are applied. In these
bounds, each prescribed Cartesian derivative reduces the power of \NSi{NS-F-P058-012}{q} by an amount that does not
depend on the truncation order.
\end{EESegment}

\begin{EESegment}{EE-TGT-P058-003}
\textbf{Step 1. Bound the derivatives of the cutoffs and their curls.}
The chain rule writes any positive-order derivative of \NSi{NS-F-P058-013}{\chi(aq)} as a sum of terms of the form
\NSFormula{NS-F-P058-014}{display}{none}
\[
 a^k\chi^{(k)}(aq)\prod_{\nu=1}^{k}\partial_z^{a_\nu}\partial_t^{b_\nu}q,
 \qquad
 \sum_\nu a_\nu=a',
 \qquad
 \sum_\nu b_\nu=b'.
\]
On the support of each such term, \NSi{NS-F-P058-015}{1/2\leq aq\leq1}. The bounds in \eqref{eq:5.28} supply factors
\NSi{NS-F-P058-016}{q^k}, which compensate for \NSi{NS-F-P058-017}{a^k} because
\NSi{NS-F-P058-018}{aq} stays bounded on the cutoff support. It follows that
\NSFormula{NS-F-P058-019}{display}{none}
\[
 |\partial_z^{a'}\partial_t^{b'}\chi(aq)|
 \leq C_{a',b'}q^{-a'D-b'},
\]
with constants that do not depend on \NSi{NS-F-P058-020}{a}.
\end{EESegment}

\begin{EESegment}{EE-TGT-P058-004}
Because \NSi{NS-F-P058-021}{q=q(z,t)}, both
\NSi{NS-F-P058-022}{\partial_{x_1}q=\partial_{x_2}q=0}. In particular, the product rule gives
\NSFormula{NS-F-P058-023}{display}{none}
\[
 \operatorname{curl}(\chi(aq)A_j)
 =\chi(aq)\operatorname{curl}A_j+a\chi'(aq)\nabla q\times A_j
\]
The second term is the commutator created by differentiating the cutoff, and the assumed derivative bounds for the potentials control it. The product rule now bounds all derivatives through
order \NSi{NS-F-P058-024}{m} of the \NSi{NS-F-P058-025}{j}\,th summand in \eqref{eq:5.34} by
\NSi{NS-F-P058-026}{\widehat C_{j,m}q^{g_j-\ell_m'}\Lambda_{\log}(q)^{\widehat P_{j,m}}}. The loss
\NSi{NS-F-P058-027}{\ell_m'} depends only on \NSi{NS-F-P058-028}{m} and
\NSi{NS-F-P058-029}{\ell_{m+1}}.
\end{EESegment}

\begin{EESegment}{EE-TGT-P058-005}
\textbf{Step 2. Choose the cutoffs and estimate the tail.}
Choose the scales \NSi{NS-F-P058-030}{a_j} one at a time so that
\NSi{NS-F-P058-031}{a_{j+1}\geq2a_j}, every cutoff \NSi{NS-F-P058-032}{\chi(a_jq)} is supported inside
\NSi{NS-F-P058-033}{q<q_0}, and
\NSFormula{NS-F-P058-034}{display}{5.37}
\begin{equation}
 \widehat C_{j,m}\Lambda_{\log}(q)^{\widehat P_{j,m}}q^{g_j/2}
 \leq2^{-j}
 \qquad(0<q\leq a_j^{-1},\ 0\leq m\leq j).
 \tag{5.37}\label{eq:5.37}
\end{equation}
At step \NSi{NS-F-P058-035}{j}, there are only finitely many requirements, and each one holds for all sufficiently small
\NSi{NS-F-P058-036}{q} because \NSi{NS-F-P058-037}{g_j>0}. All derivatives through order
\NSi{NS-F-P058-038}{m} of the resulting cutoff increment are therefore at most
\NSi{NS-F-P058-039}{2^{-j}q^{g_j/2-\ell_m'}} whenever \NSi{NS-F-P058-040}{m\leq j}.
\end{EESegment}

\begin{EESegment}{EE-TGT-P058-006}
On a compact subset where \NSi{NS-F-P058-041}{q\geq\delta>0}, every term with
\NSi{NS-F-P058-042}{a_j>\delta^{-1}} is zero, so the sum is locally finite. The curl of every vector field is
divergence-free. The field \NSi{NS-F-P058-043}{\chi(a_jq)B_j} is also divergence-free because its angular amplitude does
not depend on \NSi{NS-F-P058-044}{\theta}. Smoothness and the stated extensions across the support boundaries now follow
one term at a time.

Fix \NSi{NS-F-P058-045}{m} and \NSi{NS-F-P058-046}{J\geq\max(1,m)}. When
\NSi{NS-F-P058-047}{q<(2a_J)^{-1}}, the first \NSi{NS-F-P058-048}{J} cutoffs all equal one. Since
\NSi{NS-F-P058-049}{q<1} and \NSi{NS-F-P058-050}{g_j} does not decrease as its index rises, the remaining terms satisfy
\NSFormula{NS-F-P058-051}{display}{none}
\[
 \sum_{j>J}2^{-j}q^{g_j/2-\ell_m'}
 \leq2^{-J}q^{g_{J+1}/2-\ell_m'},
\]
which proves \eqref{eq:5.35}.
\end{EESegment}

\NSPage{059}%
\begin{EESegment}{EE-TGT-P059-001}
\textbf{Step 3. Compare the residual with one finite partial sum.}
The cutoff sum \NSi{NS-F-P059-001}{U} is now smooth, and its velocity is divergence-free. To show that
\NSi{NS-F-P059-002}{\mathcal F(U)} is flat, compare it with \NSi{NS-F-P059-003}{\mathcal F(U^{[J]})} for a finite
partial sum whose residual has sufficiently high decay. Then control the change caused by
\NSi{NS-F-P059-004}{U-U^{[J]}}. The assumed bound on \NSi{NS-F-P059-005}{|U_0|_m} and \eqref{eq:5.30} give
\NSFormula{NS-F-P059-006}{display}{5.38}
\begin{equation}
 |U^{[J]}|_m\leq C_{J,m}q^{-K_m}\Lambda_{\log}(q)^{P_{J,m}},
 \tag{5.38}\label{eq:5.38}
\end{equation}
where \NSi{NS-F-P059-007}{K_m\geq0} does not depend on \NSi{NS-F-P059-008}{J}. This is because every increment has
\NSi{NS-F-P059-009}{g_j>0}, and taking the curl uses one additional derivative of the potential.
\end{EESegment}

\begin{EESegment}{EE-TGT-P059-002}
The bounds \NSi{NS-F-P059-010}{s,d} on order and degree are the ones from \eqref{eq:5.31}. Subtract the two products
one factor at a time, then use the product rule and \eqref{eq:5.32}. For
\NSi{NS-F-P059-011}{e=U-U^{[J]}}, this gives
\NSFormula{NS-F-P059-012}{display}{5.39}
\begin{equation}
 |\mathcal F(U)-\mathcal F(U^{[J]})|_m
 \leq C_mq^{-H_m}|e|_{m+s}
 \bigl(1+|U^{[J]}|_{m+s}+|e|_{m+s}\bigr)^{d-1},
 \tag{5.39}\label{eq:5.39}
\end{equation}
where \NSi{NS-F-P059-013}{H_m} does not depend on \NSi{NS-F-P059-014}{J}. The finitely many logarithmic factors from
the coefficients are absorbed into this fixed power of \NSi{NS-F-P059-015}{q^{-1}}. A polynomial that does not depend
on \NSi{NS-F-P059-016}{U} has zero difference, so it needs no estimate here.
\end{EESegment}

\begin{EESegment}{EE-TGT-P059-003}
Fix a derivative order \NSi{NS-F-P059-017}{m} and a decay order \NSi{NS-F-P059-018}{N}. Choose
\NSi{NS-F-P059-019}{J\geq m+s} large enough that
\NSFormula{NS-F-P059-020}{display}{none}
\[
 \frac{g_{J+1}}2-\ell_{m+s}'
 \geq N+H_m+(d-1)(K_{m+s}+1),
 \qquad
 \rho_J-K_m^{\mathcal F}\geq N+1.
\]
Also choose \NSi{NS-F-P059-021}{J} large enough that
\NSi{NS-F-P059-022}{g_{J+1}/2-\ell_{m+s}'} is nonnegative, and then keep
\NSi{NS-F-P059-023}{J} fixed. Shrink the \NSi{NS-F-P059-024}{q}-neighborhood until
\NSi{NS-F-P059-025}{\chi(a_jq)=1} for \NSi{NS-F-P059-026}{j\leq J}. In this neighborhood, one power of
\NSi{NS-F-P059-027}{q^{-1}} absorbs the constants and logarithms in \eqref{eq:5.38}, and \eqref{eq:5.35} gives
\NSi{NS-F-P059-028}{|e|_{m+s}\leq1}. Formula~\eqref{eq:5.39} is then
\NSi{NS-F-P059-029}{O(q^N)}.
\end{EESegment}

\begin{EESegment}{EE-TGT-P059-004}
The same conclusion holds for \eqref{eq:5.33}, after one inverse power of the small parameter absorbs its logarithm and the
flatness of the fixed \NSi{NS-F-P059-030}{E_{J,m}} is used. This proves \eqref{eq:5.36}. The choices are made in this
order. First choose \NSi{NS-F-P059-031}{(m,N)}, then \NSi{NS-F-P059-032}{J}, and then the neighborhood and its
constants. The comparison uses only the remainder \NSi{NS-F-P059-033}{E_{J,m}} at that fixed truncation.
\end{EESegment}
\end{proof}

\begin{EESegment}{EE-TGT-P059-005}
A finite initial block whose decay orders are not positive can be given one common cutoff, applied as in
\eqref{eq:5.34}, and included in the base \NSi{NS-F-P059-034}{U_0}. That cutoff is one near
\NSi{NS-F-P059-035}{q=0}, so every comparison with a finite partial sum is unchanged there.
\end{EESegment}

\begin{EESegment}{EE-TGT-P059-006}
For smooth vector fields \NSi{NS-F-P059-036}{v,e} and smooth scalar fields \NSi{NS-F-P059-037}{p,\pi}, the physical
residual changes by exactly
\NSFormula{NS-F-P059-038}{display}{none}
\[
 \mathscr R(v+e,p+\pi)-\mathscr R(v,p)
 =\partial_te-\Delta e+\nabla\pi+(v\mathbin\cdot\nabla)e+(e\mathbin\cdot\nabla)v
 +(e\mathbin\cdot\nabla)e.
\]
For \NSi{NS-F-P059-039}{\mathcal F_{\mathrm{slow}}}, also add the change in the linear tangential stress-divergence
term from \eqref{eq:5.24}. On the common annular stress support, the cylindrical coefficients and every prescribed
Cartesian derivative of them are bounded by a fixed negative power of \NSi{NS-F-P059-040}{q}. In the final correction
sequence, \NSi{NS-F-P059-041}{\mathcal F=\mathscr R}, and there is no stress variable.
\end{EESegment}

\begin{EESegment}{EE-TGT-P059-007}
The same cutoff choices can satisfy further normalized estimates once those estimates have been proved for the
coefficients. More exactly, at step \NSi{NS-F-P059-042}{j}, any finite list of conditions
\NSFormula{NS-F-P059-043}{display}{5.40}
\begin{equation}
 B_{j,k}\Lambda_{\log}(q)^{P_{j,k}}q^{\gamma_{j,k}}\leq2^{-j},
 \qquad \gamma_{j,k}>0,
 \tag{5.40}\label{eq:5.40}
\end{equation}
can be added to \eqref{eq:5.37}. For derivatives of the normalized expansion coefficients and boundary weights, impose
these extra coefficient estimates through the corresponding product rules for the cutoffs. To recover a full expansion
after a fixed truncation \NSi{NS-F-P059-044}{N} and at a fixed derivative order \NSi{NS-F-P059-045}{m}, choose
\NSi{NS-F-P059-046}{M\geq\max(N,m)} with \NSi{NS-F-P059-047}{g_{M+1}/2\geq g_{N+1}}. The finite block from
\NSi{NS-F-P059-048}{N+1} through \NSi{NS-F-P059-049}{M} keeps its original orders near
\NSi{NS-F-P059-050}{q=0}. The bound \eqref{eq:5.35} controls the rest of the tail at order
\NSi{NS-F-P059-051}{g_{N+1}}, with the fixed exponent reduction \NSi{NS-F-P059-052}{\ell_m'}.
\end{EESegment}

\NSPage{060}%
\begin{EESegment}{EE-TGT-P060-001}
\subsection{The base field after summation}\label{sec:5.5}
The finite tuples \NSi{NS-F-P060-001}{\mathcal U^{[N]}} from Proposition~\ref{prop:5.3} now have the residual estimates
needed by Lemma~\ref{lem:5.4}. Apply that lemma to the coefficient fields and their potentials
\NSi{NS-F-P060-002}{\mathcal A_n} from \eqref{eq:5.27} to build the smooth background
\NSi{NS-F-P060-003}{(u_B,p_B)}. The construction must also keep the leading profile as in \eqref{eq:5.42}, the exact heat
field outside the active region from \eqref{eq:4.29}, and the weighted stress bounds in \eqref{eq:5.43}. Those bounds are
used to choose the wave amplitudes and their corrections in Propositions~\ref{prop:7.5} and~\ref{prop:7.6}.
\end{EESegment}

\begin{proposition}\label{prop:5.5}
\begin{EESegment}{EE-TGT-P060-002}
There are smooth axisymmetric fields \NSi{NS-F-P060-004}{(u_B,p_B)} for \NSi{NS-F-P060-005}{q>0}, together with two
physical tangential stress components \NSi{NS-F-P060-006}{\mathcal T_{\mathrm{phys}}} supported where
\NSi{NS-F-P060-007}{X_a\leq X(r,z,t)\leq X_b}, such that
\NSi{NS-F-P060-008}{\operatorname{div}u_B=0} and
\NSFormula{NS-F-P060-009}{display}{5.41}
\begin{equation}
 \mathscr R(u_B,p_B)
 =-(\partial_r+2/r)\mathcal T_{\mathrm{phys},\theta}e_\theta
  -(\partial_r+1/r)\mathcal T_{\mathrm{phys},z}e_z+\mathcal E_B.
 \tag{5.41}\label{eq:5.41}
\end{equation}
Fix \NSi{NS-F-P060-010}{0<X_{\max}<\infty}. Uniformly on the profile rectangle
\NSi{NS-F-P060-011}{[0,X_{\max}]\times[-1,1]} in \NSi{NS-F-P060-012}{(X,\eta)}, every Cartesian space--time
derivative \NSi{NS-F-P060-013}{\partial^\alpha} and every \NSi{NS-F-P060-014}{M>0} satisfy
\NSi{NS-F-P060-015}{|\partial^\alpha\mathcal E_B|\leq C_{\alpha,M}q^M} as
\NSi{NS-F-P060-016}{q\downarrow0}. The constants can depend on \NSi{NS-F-P060-017}{X_{\max}}. The formal expansion of
\NSi{NS-F-P060-018}{(u_B,p_B,\mathcal T_{\mathrm{phys}})} has the finite truncations
\NSi{NS-F-P060-019}{\mathcal U^{[N]}} from Proposition~\ref{prop:5.3}, with the physical Stokes streamfunctions
\NSi{NS-F-P060-020}{\mathcal S_n} from \eqref{eq:5.27} summed before differentiation.
\end{EESegment}

\begin{EESegment}{EE-TGT-P060-003}
Fix radial endpoints \NSi{NS-F-P060-021}{0<X_{\mathrm{lo}}<X_a<X_b<X_{\mathrm{hi}}<\infty}. The rectangle
\NSi{NS-F-P060-022}{[X_{\mathrm{lo}},X_{\mathrm{hi}}]\times[-1,1]} is a fixed radial enlargement of the closed active
annulus \NSi{NS-F-P060-023}{[X_a,X_b]\times[-1,1]} from Item~\textup{(ii)} of Theorem~\ref{thm:4.6}, and
\NSi{NS-F-P060-024}{\eta} runs through its full closed parameter interval. On this enlarged rectangle, take any fixed
composition \NSi{NS-F-P060-025}{\mathscr D^I} of
\NSi{NS-F-P060-026}{q\partial_q,\partial_X,\partial_\eta}. Then, as \NSi{NS-F-P060-027}{q\downarrow0},
\NSFormula{NS-F-P060-028}{display}{5.42}
\begin{equation}
 \begin{aligned}
  \mathscr D^I(q^Au_{\theta,B}-E_0)&=O_I(q^{2h}),
  &\mathscr D^I(q^Au_{z,B}-U_0)&=O_I(q^{2h}),\\
  \mathscr D^I(q^{1/2}u_{r,B}-V_0/\sqrt{2X})&=O_I(q^{2h}),
  &\mathscr D^I(q^Au_{r,B})&=O_I(q^h).
 \end{aligned}
 \tag{5.42}\label{eq:5.42}
\end{equation}
The constants in these estimates can depend on the fixed radial endpoints. With
\NSi{NS-F-P060-029}{\widehat{\mathcal T}=q^{A+1/2}\mathcal T_{\mathrm{phys}}}, the weighted comparison uses
\NSi{NS-F-P060-030}{\zeta,\delta} from Item~\textup{(iv)} of Theorem~\ref{thm:4.6} and holds whenever
\NSi{NS-F-P060-031}{X_a<X<X_b} and \NSi{NS-F-P060-032}{-1\leq\eta\leq1}.
\NSFormula{NS-F-P060-033}{display}{5.43}
\begin{equation}
 |\mathscr D^I(\widehat{\mathcal T}-\mathcal T_0)|
 \leq C_Iq^{2h}\zeta\delta^{-N_I}
 \leq C_Iq^h\zeta\delta^{-N_I}
 \qquad(0<q\leq1).
 \tag{5.43}\label{eq:5.43}
\end{equation}
\end{EESegment}

\begin{EESegment}{EE-TGT-P060-004}
Every fixed physical space--time derivative of
\NSi{NS-F-P060-034}{u_B,p_B,\mathcal T_{\mathrm{phys}},\mathcal E_B} extends continuously to
\NSi{NS-F-P060-035}{\eta=\pm1} on compact sets where \NSi{NS-F-P060-036}{q>0}. The positive-order corrections
\NSi{NS-F-P060-037}{u_B-u^{(0)}} and \NSi{NS-F-P060-038}{p_B-p^{(0)}} are zero whenever
\NSi{NS-F-P060-039}{X\geq X_+}, with \NSi{NS-F-P060-040}{X_+} from Lemma~\ref{lem:5.2}. In particular, the exact
exterior of the leading profile from \eqref{eq:4.29} stays unchanged. On the reserved interval
\NSi{NS-F-P060-041}{X\in\mathcal I_{\mathrm{mean}}\subset(X_a,X_v)} from Item~\textup{(vi)} of
Theorem~\ref{thm:4.6}, for every \NSi{NS-F-P060-042}{-1\leq\eta\leq1}, the base velocity agrees exactly with the leading
profile.
\NSFormula{NS-F-P060-043}{display}{5.44}
\begin{equation}
 u_B=u^{(0)},
 \qquad u_{z,B}=0,
 \qquad q^Au_{\theta,B}=E_0=c_{\mathrm{patch}}(1+\eta^2)^{-1}X^{-1/2-\lambda}.
 \tag{5.44}\label{eq:5.44}
\end{equation}
Here \NSi{NS-F-P060-044}{u^{(0)}} is the velocity determined by the leading profile in \eqref{eq:4.3}.
\end{EESegment}
\end{proposition}

\begin{proof}
\begin{EESegment}{EE-TGT-P060-005}
\textbf{Step 1. Check the assumptions of the summation lemma and choose the cutoffs.}
Use the potentials in \eqref{eq:5.27} and the physical tuples listed before Lemma~\ref{lem:5.4}, while leaving order zero
unchanged. The Cartesian formula for \NSi{NS-F-P060-045}{\mathcal A_n} proves regularity at the axis, and
Lemma~\ref{lem:5.2} gives the common supports and all fixed derivative bounds for the coefficients. By \eqref{eq:5.26},
because \NSi{NS-F-P060-046}{D<1/2}, the reduction in the power of \NSi{NS-F-P060-047}{q} for the coefficient tuple in
\eqref{eq:5.27} and \eqref{eq:5.1}, including the stress components, is bounded by
\NSi{NS-F-P060-048}{\ell_m=2A+m}. The coefficient constants
absorb the powers of \NSi{NS-F-P060-049}{n} produced by differentiating
\NSi{NS-F-P060-050}{q^{\lambda_n}}. For the cutoffs, set \NSi{NS-F-P060-051}{\sigma=c_nq}. Then every integer
\NSi{NS-F-P060-052}{j\geq0} satisfies
\NSFormula{NS-F-P060-053}{display}{none}
\[
 (q\partial_q)^j\chi(c_nq)
 =(\sigma\partial_\sigma)^j\chi(\sigma)\big|_{\sigma=c_nq},
\]
\NSFormula{NS-F-P060-054}{display}{none}
\[
 \sup_{n\geq1,\,q>0}|(q\partial_q)^j\chi(c_nq)|
 \leq\sup_{\sigma>0}|(\sigma\partial_\sigma)^j\chi(\sigma)|
 =:C_j<\infty.
\]
\end{EESegment}
